\documentclass[tikz,border=4pt]{standalone}
% ==========================================================================
% tikz-tcee.tex — Estilos comunes de los gráficos TikZ del temario
% Autor: Víctor Gutiérrez Marcos
%
% Cada gráfico es un fichero standalone en fuentes/<tema>/graficos/:
%
%   \documentclass[tikz,border=4pt]{standalone}
%   % ==========================================================================
% tikz-tcee.tex — Estilos comunes de los gráficos TikZ del temario
% Autor: Víctor Gutiérrez Marcos
%
% Cada gráfico es un fichero standalone en fuentes/<tema>/graficos/:
%
%   \documentclass[tikz,border=4pt]{standalone}
%   % ==========================================================================
% tikz-tcee.tex — Estilos comunes de los gráficos TikZ del temario
% Autor: Víctor Gutiérrez Marcos
%
% Cada gráfico es un fichero standalone en fuentes/<tema>/graficos/:
%
%   \documentclass[tikz,border=4pt]{standalone}
%   \input{../../../latex/tikz-tcee}
%   \begin{document}
%   \begin{tikzpicture}[tcee]
%     \ejes{Q}{P}{6}{5}
%     \draw[curva1] (0.5,4.5) -- (5,0.8) node[etiqueta,right] {$D$};
%     \capa{2}{\draw[curva2] ...;}   % en el vídeo aparece en el paso 2
%   \end{tikzpicture}
%   \end{document}
%
% Paleta: la de la web (morado, dorado, salvia) más tres colores de apoyo
% distinguibles en blanco y negro por el trazo.
% ==========================================================================
\usepackage[T1]{fontenc}
\usepackage{newpxtext,newpxmath}
\usepackage{xcolor}
\usepackage{tikz,pgfplots}
\pgfplotsset{compat=1.18}
\usetikzlibrary{arrows.meta,calc,positioning,patterns,decorations.pathreplacing,intersections,shapes.misc,backgrounds}

\definecolor{tceemorado}{HTML}{5F2987}
\definecolor{tceedorado}{HTML}{B8860B}
\definecolor{tceeverde}{HTML}{3C7D3A}
\definecolor{tceeazul}{HTML}{1F5F99}
\definecolor{tceerojo}{HTML}{B22222}
\definecolor{tceegris}{HTML}{767171}
\definecolor{tceesalvia}{HTML}{E2EFD9}

% Capas para el vídeo: \capa{n}{…} dibuja su contenido solo a partir del paso n.
% En el PDF, la web y Word se ve todo; generar-video.py compila el gráfico una
% vez por paso con \def\tceepaso{k} para que cada elemento aparezca al narrarlo.
\providecommand{\tceepaso}{1000}
\newcommand{\capa}[2]{\ifnum#1>\tceepaso\relax\else#2\fi}

\tikzset{
  tcee/.style={line cap=round,line join=round,font=\small,>={Stealth[length=2.2mm]}},
  eje/.style={->,thick,black},
  curva1/.style={very thick,tceemorado},
  curva2/.style={very thick,tceeazul},
  curva3/.style={very thick,tceeverde},
  curva4/.style={very thick,tceedorado},
  curvanueva/.style={very thick,tceerojo},
  desplazada/.style={very thick,dashed},
  guia/.style={thin,dashed,tceegris},
  punto/.style={circle,fill=black,inner sep=1.3pt},
  etiqueta/.style={font=\small,inner sep=2pt},
  flecha/.style={->,thick,tceerojo},
  area/.style={fill=tceemorado,fill opacity=0.18,draw=none},
  area2/.style={fill=tceedorado,fill opacity=0.22,draw=none},
}

% \ejes{etiqueta x}{etiqueta y}{largo x}{alto y}
\newcommand{\ejes}[4]{%
  \draw[eje] (0,0) -- (#3,0) node[below left=2pt and 0pt] {#1};
  \draw[eje] (0,0) -- (0,#4) node[left] {#2};
  \node[below left] at (0,0) {$0$};}
% \proyeccion{(x,y)}{etq x}{etq y}: líneas guía a los ejes con etiquetas
\newcommand{\proyeccion}[3]{%
  \draw[guia] let \p1=#1 in (\x1,0) node[below] {#2} |- (0,\y1) node[left] {#3};}

%   \begin{document}
%   \begin{tikzpicture}[tcee]
%     \ejes{Q}{P}{6}{5}
%     \draw[curva1] (0.5,4.5) -- (5,0.8) node[etiqueta,right] {$D$};
%     \capa{2}{\draw[curva2] ...;}   % en el vídeo aparece en el paso 2
%   \end{tikzpicture}
%   \end{document}
%
% Paleta: la de la web (morado, dorado, salvia) más tres colores de apoyo
% distinguibles en blanco y negro por el trazo.
% ==========================================================================
\usepackage[T1]{fontenc}
\usepackage{newpxtext,newpxmath}
\usepackage{xcolor}
\usepackage{tikz,pgfplots}
\pgfplotsset{compat=1.18}
\usetikzlibrary{arrows.meta,calc,positioning,patterns,decorations.pathreplacing,intersections,shapes.misc,backgrounds}

\definecolor{tceemorado}{HTML}{5F2987}
\definecolor{tceedorado}{HTML}{B8860B}
\definecolor{tceeverde}{HTML}{3C7D3A}
\definecolor{tceeazul}{HTML}{1F5F99}
\definecolor{tceerojo}{HTML}{B22222}
\definecolor{tceegris}{HTML}{767171}
\definecolor{tceesalvia}{HTML}{E2EFD9}

% Capas para el vídeo: \capa{n}{…} dibuja su contenido solo a partir del paso n.
% En el PDF, la web y Word se ve todo; generar-video.py compila el gráfico una
% vez por paso con \def\tceepaso{k} para que cada elemento aparezca al narrarlo.
\providecommand{\tceepaso}{1000}
\newcommand{\capa}[2]{\ifnum#1>\tceepaso\relax\else#2\fi}

\tikzset{
  tcee/.style={line cap=round,line join=round,font=\small,>={Stealth[length=2.2mm]}},
  eje/.style={->,thick,black},
  curva1/.style={very thick,tceemorado},
  curva2/.style={very thick,tceeazul},
  curva3/.style={very thick,tceeverde},
  curva4/.style={very thick,tceedorado},
  curvanueva/.style={very thick,tceerojo},
  desplazada/.style={very thick,dashed},
  guia/.style={thin,dashed,tceegris},
  punto/.style={circle,fill=black,inner sep=1.3pt},
  etiqueta/.style={font=\small,inner sep=2pt},
  flecha/.style={->,thick,tceerojo},
  area/.style={fill=tceemorado,fill opacity=0.18,draw=none},
  area2/.style={fill=tceedorado,fill opacity=0.22,draw=none},
}

% \ejes{etiqueta x}{etiqueta y}{largo x}{alto y}
\newcommand{\ejes}[4]{%
  \draw[eje] (0,0) -- (#3,0) node[below left=2pt and 0pt] {#1};
  \draw[eje] (0,0) -- (0,#4) node[left] {#2};
  \node[below left] at (0,0) {$0$};}
% \proyeccion{(x,y)}{etq x}{etq y}: líneas guía a los ejes con etiquetas
\newcommand{\proyeccion}[3]{%
  \draw[guia] let \p1=#1 in (\x1,0) node[below] {#2} |- (0,\y1) node[left] {#3};}

%   \begin{document}
%   \begin{tikzpicture}[tcee]
%     \ejes{Q}{P}{6}{5}
%     \draw[curva1] (0.5,4.5) -- (5,0.8) node[etiqueta,right] {$D$};
%     \capa{2}{\draw[curva2] ...;}   % en el vídeo aparece en el paso 2
%   \end{tikzpicture}
%   \end{document}
%
% Paleta: la de la web (morado, dorado, salvia) más tres colores de apoyo
% distinguibles en blanco y negro por el trazo.
% ==========================================================================
\usepackage[T1]{fontenc}
\usepackage{newpxtext,newpxmath}
\usepackage{xcolor}
\usepackage{tikz,pgfplots}
\pgfplotsset{compat=1.18}
\usetikzlibrary{arrows.meta,calc,positioning,patterns,decorations.pathreplacing,intersections,shapes.misc,backgrounds}

\definecolor{tceemorado}{HTML}{5F2987}
\definecolor{tceedorado}{HTML}{B8860B}
\definecolor{tceeverde}{HTML}{3C7D3A}
\definecolor{tceeazul}{HTML}{1F5F99}
\definecolor{tceerojo}{HTML}{B22222}
\definecolor{tceegris}{HTML}{767171}
\definecolor{tceesalvia}{HTML}{E2EFD9}

% Capas para el vídeo: \capa{n}{…} dibuja su contenido solo a partir del paso n.
% En el PDF, la web y Word se ve todo; generar-video.py compila el gráfico una
% vez por paso con \def\tceepaso{k} para que cada elemento aparezca al narrarlo.
\providecommand{\tceepaso}{1000}
\newcommand{\capa}[2]{\ifnum#1>\tceepaso\relax\else#2\fi}

\tikzset{
  tcee/.style={line cap=round,line join=round,font=\small,>={Stealth[length=2.2mm]}},
  eje/.style={->,thick,black},
  curva1/.style={very thick,tceemorado},
  curva2/.style={very thick,tceeazul},
  curva3/.style={very thick,tceeverde},
  curva4/.style={very thick,tceedorado},
  curvanueva/.style={very thick,tceerojo},
  desplazada/.style={very thick,dashed},
  guia/.style={thin,dashed,tceegris},
  punto/.style={circle,fill=black,inner sep=1.3pt},
  etiqueta/.style={font=\small,inner sep=2pt},
  flecha/.style={->,thick,tceerojo},
  area/.style={fill=tceemorado,fill opacity=0.18,draw=none},
  area2/.style={fill=tceedorado,fill opacity=0.22,draw=none},
}

% \ejes{etiqueta x}{etiqueta y}{largo x}{alto y}
\newcommand{\ejes}[4]{%
  \draw[eje] (0,0) -- (#3,0) node[below left=2pt and 0pt] {#1};
  \draw[eje] (0,0) -- (0,#4) node[left] {#2};
  \node[below left] at (0,0) {$0$};}
% \proyeccion{(x,y)}{etq x}{etq y}: líneas guía a los ejes con etiquetas
\newcommand{\proyeccion}[3]{%
  \draw[guia] let \p1=#1 in (\x1,0) node[below] {#2} |- (0,\y1) node[left] {#3};}

\usepackage{pgfplotstable}
\begin{document}
% Ley de Engel en España: % del gasto en alimentos y bebidas no alcohólicas
% por quintil de gasto, EPF 2025 (INE, tabla 73828). Datos en ../datos/.
\pgfplotstableread[col sep=comma,comment chars={\#}]{../datos/epf2025-alimentos-quintil.csv}\epf
\pgfplotstablegetelem{5}{porcentaje_alimentos}\of\epf
\let\epftotal\pgfplotsretval
\begin{tikzpicture}[tcee]
  \begin{axis}[
      width=9cm, height=6cm,
      ybar, bar width=0.9cm,
      ymin=0, ymax=25, ytick={0,5,...,25},
      yticklabel style={/pgf/number format/precision=0, /pgf/number format/fixed zerofill=false},
      xmin=-0.6, xmax=4.6, xtick={0,...,4},
      xticklabels={{Q1\\\footnotesize(menos gasto)},Q2,Q3,Q4,{Q5\\\footnotesize(más gasto)}},
      xticklabel style={align=center},
      xlabel={Quintil de gasto de los hogares},
      ylabel={\% del gasto total},
      clip=false, axis lines*=left, tick align=outside,
      ymajorgrids, grid style={tceegris!25},
      /pgf/number format/.cd, use comma, fixed, fixed zerofill, precision=2,
    ]
    \capa{1}{
      \addplot[fill=tceemorado!80, draw=tceemorado,
        skip coords between index={5}{6},
        nodes near coords, nodes near coords style={font=\footnotesize,black}]
        table[x expr=\coordindex, y=porcentaje_alimentos] {\epf};
    }
    \capa{2}{
      \draw[very thick, dashed, tceedorado]
        (axis cs:-0.6,\epftotal) -- (axis cs:4.6,\epftotal)
        node[pos=1, right, align=left, font=\footnotesize, text=tceedorado!70!black]
        {Total\\\pgfmathprintnumber[use comma,fixed,fixed zerofill,precision=2]{\epftotal}\,\%};
    }
  \end{axis}
\end{tikzpicture}
\end{document}
